% ============================================================================
%  Vigil: Explainable Concept Drift Detection in Network Traffic Streams
%  via Feature-Level Reconstruction Error Attribution
%
%  arXiv preprint — cs.LG / cs.CR
%  Author: Venkateswara Sahu
% ============================================================================
\documentclass[10pt,conference]{IEEEtran}

\usepackage{cite}
\usepackage{amsmath,amssymb}
\usepackage{graphicx}
\usepackage{booktabs}
\usepackage{hyperref}
\usepackage{url}
\usepackage{xcolor}
\usepackage{algorithm}
\usepackage{algorithmic}
\usepackage{microtype}

\hypersetup{
  colorlinks=true,
  linkcolor=blue!70!black,
  citecolor=green!50!black,
  urlcolor=blue!70!black
}

\begin{document}

% ── Title ─────────────────────────────────────────────────────────────────────
\title{Vigil: Explainable Concept Drift Detection\\
       in Network Traffic Streams via Feature-Level\\
       Reconstruction Error Attribution}

\author{
  \IEEEauthorblockN{Venkateswara Sahu}
  \IEEEauthorblockA{
    \textit{B.Tech (Hons.) CSE, Data Science \& Data Engineering}\\
    \textit{Lovely Professional University}\\
    Punjab, India\\
    \href{mailto:venkateswarsahu000@gmail.com}{venkateswarsahu000@gmail.com}
  }
}

\maketitle

% ── Abstract ──────────────────────────────────────────────────────────────────
\begin{abstract}
Detecting concept drift in high-dimensional, non-stationary network
traffic streams is critical for maintaining the reliability of
production machine learning systems in cybersecurity.
Existing drift detectors identify \emph{when} a distribution shift
occurs but provide no guidance on \emph{which} input features caused
it---a significant limitation for network operators and security analysts
who must act on alerts.
We present \textbf{Vigil}, an open-source Python library that addresses this
explainability gap through \emph{DriftAttributor}: a per-feature
reconstruction-error attribution module built on top of dual,
mirrored autoencoders (A and A\textsubscript{KC}) and a replicated
Welch T-test.
When drift is confirmed, DriftAttributor ranks input features by
their relative increase in reconstruction error between the stable
reference distribution and the drifted stream, surfacing the exact
network features responsible.
Evaluated on the NSL-KDD network intrusion benchmark
(125,973 samples, 122 features, 22 attack classes),
Vigil achieves \textbf{93.3\% drift detection precision} and
\textbf{100\% novel-class recall}, detecting shifts within
\textbf{one chunk (200 packets)} of injection---
outperforming classical univariate detectors (ADWIN, KSWIN,
Page-Hinkley) that achieve near-zero recall on the same benchmark.
Vigil is available at
\url{https://pypi.org/project/vigil-drift/} and
\url{https://github.com/Venkateswara-Sahu/OWADD}.
\end{abstract}

\begin{IEEEkeywords}
concept drift, network intrusion detection, autoencoders,
feature attribution, unsupervised learning, streaming ML
\end{IEEEkeywords}

% ── 1. Introduction ───────────────────────────────────────────────────────────
\section{Introduction}

Modern network security systems rely on machine learning models trained
on historical traffic data. In production, the statistical properties of
incoming traffic evolve continuously: new attack campaigns emerge, protocol
usage shifts, and infrastructure changes alter baseline behaviour.
This phenomenon---\emph{concept drift}---degrades model performance
silently and unpredictably~\cite{gama2014survey}.

Detecting drift in network streams is challenging for three reasons.
First, network traffic is \emph{high-dimensional}: modern feature
engineering pipelines for intrusion detection produce 100+ features
per connection record~\cite{tavallaee2009nsl}.
Second, labels are unavailable in real time: a Security Operations Centre
(SOC) cannot wait for ground-truth attack labels to determine whether
the traffic distribution has changed.
Third, existing detectors provide \emph{no explainability}: when an alert
fires, analysts do not know which features triggered it, making
triage slow and error-prone.

Classical drift detectors such as ADWIN~\cite{bifet2007adwin},
KSWIN~\cite{raab2020kswin}, and the Page-Hinkley
test~\cite{page1954continuous} operate on a single scalar error signal.
Applying them to high-dimensional traffic data requires reducing the
feature space to a univariate proxy (e.g., classifier error rate),
which loses distributional information and results in poor recall, as
we demonstrate empirically in Section~\ref{sec:experiments}.

We present \textbf{Vigil}, which makes two contributions:
\begin{enumerate}
  \item A \emph{production-ready, open-source implementation} of the
        dual-autoencoder drift detection framework of
        Huang~et~al.~\cite{huang2025owadd}, extending it with
        a Kafka-native streaming pipeline, Airflow auto-retrain DAGs,
        MLflow experiment tracking, and an 81\% test-coverage CI/CD suite.
  \item \textbf{DriftAttributor}: a novel, label-free module that
        computes per-feature reconstruction-error deltas when drift is
        detected, ranking features by their causal contribution to the
        distribution shift. To our knowledge, no existing open-source
        drift detector provides feature-level attribution without labels.
\end{enumerate}

% ── 2. Related Work ───────────────────────────────────────────────────────────
\section{Related Work}

\subsection{Classical Drift Detection}
ADWIN~\cite{bifet2007adwin} maintains an adaptive sliding window and
detects drift when two sub-windows have statistically different means.
KSWIN~\cite{raab2020kswin} applies the Kolmogorov-Smirnov test over a
fixed-size window. The Page-Hinkley test~\cite{page1954continuous}
monitors cumulative deviation from an observed mean.
All three are univariate and require a scalar error signal,
making them poorly suited for high-dimensional feature spaces.

\subsection{Autoencoder-Based Anomaly and Drift Detection}
Autoencoders have been widely used for anomaly detection in network
traffic~\cite{sakurada2014anomaly}, where high reconstruction error
indicates out-of-distribution samples. Huang~et~al.~\cite{huang2025owadd}
propose a dual-autoencoder framework: an adaptive autoencoder (A)
that tracks the current distribution, and a frozen mirror autoencoder
(A\textsubscript{KC}) that retains the initial known-class distribution
for novelty detection. We build Vigil on top of this architecture
and extend it with feature-level attribution.

\subsection{Explainable Drift Detection}
SHAP~\cite{lundberg2017shap} and LIME~\cite{ribeiro2016lime}
explain individual model predictions but do not operate on
streaming distributional shifts.
Explainable drift detection is an emerging area:
Hinder~et~al.~\cite{hinder2023explain} propose post-hoc explanation
of detected drift, but require labelled reference data.
Vigil's DriftAttributor is fully unsupervised and integrated
into the detection pipeline.

% ── 3. Method ─────────────────────────────────────────────────────────────────
\section{Method}
\label{sec:method}

\subsection{Problem Formulation}
Let $\mathcal{S} = \{X_1, X_2, \ldots\}$ be a non-stationary data
stream where each chunk $X_t \in \mathbb{R}^{n \times d}$ contains
$n$ samples with $d$ features. We assume access to an initial reference
chunk $X_0$ of known-class (normal) traffic. Our goal is to detect, for
each $t > 0$, whether $P(X_t) \neq P(X_0)$, and if so, to rank all $d$
features by their contribution to the distribution shift---without labels.

\subsection{Dual Autoencoder Architecture}

\textbf{Autoencoder A} (adaptive): A standard fully-connected
autoencoder with encoder $f_\theta: \mathbb{R}^d \to \mathbb{R}^h$
and decoder $g_\phi: \mathbb{R}^h \to \mathbb{R}^d$. After each
confirmed drift event, A is fine-tuned on recent data to track the
evolving distribution.

\textbf{Autoencoder A\textsubscript{KC}} (frozen mirror): Initialised
with identical weights to A after the offline training phase,
then frozen permanently. A\textsubscript{KC} serves as a reference to
the original known-class distribution.

Both autoencoders are trained with mean-squared reconstruction error:
\begin{equation}
  \mathcal{L}(X) = \frac{1}{n} \sum_{i=1}^{n} \|x_i - g_\phi(f_\theta(x_i))\|^2
\end{equation}

\subsection{Drift Detection via Replicated T-Test}

For chunk $X_t$, let $E_A^{(t)} = \{e_A(x_i)\}_{i=1}^n$ be the
per-sample reconstruction errors from autoencoder A, and let
$E_{\mathrm{ref}}$ be the errors buffered from the reference distribution.

We perform $r = 15$ replications of Welch's two-sample T-test,
each drawing $s = 30$ samples without replacement from
$E_A^{(t)}$ and $E_{\mathrm{ref}}$.
The drift threshold $\delta$ is the minimum fraction of replications
that must reject $H_0$ (equal means) to declare drift:
\begin{equation}
  \text{drift}(X_t) = \mathbf{1}\left[
    \frac{1}{r}\sum_{k=1}^{r} \mathbf{1}[p_k < 0.05] \geq \delta
  \right]
\end{equation}
We set $\delta = 0.10$ (10\%) in all experiments.

\subsection{Novel-Class Recognition via KDE}

Autoencoder A\textsubscript{KC} is used to identify samples from
previously unseen classes. Its reconstruction errors on normal data
are used to fit a Kernel Density Estimator (KDE) with bandwidth
selected by Scott's rule~\cite{scott1992multivariate}.

For each new sample $x$, if its density $\hat{p}(e_{\mathrm{KC}}(x))$
falls below the 5\textsuperscript{th} percentile of the reference
density distribution, the sample is flagged as potentially novel.
The chunk-level novelty proportion is reported as the fraction of
flagged samples.

\subsection{DriftAttributor: Feature-Level Attribution}
\label{sec:attributor}

This is the primary novel contribution of this work.
When $\text{drift}(X_t) = 1$, DriftAttributor computes the
\emph{reconstruction error delta} for each feature $j$:
\begin{equation}
  \Delta_j = \frac{1}{n}\sum_{i=1}^n \left(x_{ij} - \hat{x}_{ij}\right)^2
           - \frac{1}{|\mathcal{B}|}\sum_{x \in \mathcal{B}}
             \left(x_j - \hat{x}_j\right)^2
\end{equation}
where $\hat{x}$ denotes the reconstructed output of autoencoder A,
and $\mathcal{B}$ is the reference buffer.
$\Delta_j$ measures how much more difficult feature $j$ has become
for the autoencoder to reconstruct under the current distribution.

Features are ranked by $\Delta_j$ in descending order and the top-$k$
are reported (default $k=5$).
The contribution percentage of feature $j$ is:
\begin{equation}
  c_j = \frac{\Delta_j}{\sum_{i=1}^d \max(\Delta_i, 0)} \times 100\%
\end{equation}

\begin{algorithm}[t]
\caption{Vigil Online Detection with DriftAttributor}
\label{alg:vigil}
\begin{algorithmic}[1]
\REQUIRE Chunk $X_t$, autoencoders A, A\textsubscript{KC},
         reference buffer $\mathcal{B}$, threshold $\delta$
\STATE $E_A \leftarrow \mathrm{ReconstructionErrors}(A, X_t)$
\STATE $E_{\mathrm{KC}} \leftarrow \mathrm{ReconstructionErrors}(A_{\mathrm{KC}}, X_t)$
\STATE $\mathrm{drift} \leftarrow \mathrm{ReplicatedTTest}(E_A, \mathcal{B}, \delta)$
\STATE $\mathrm{novelty} \leftarrow \mathrm{KDEDetect}(E_{\mathrm{KC}})$
\IF{drift}
  \STATE $\Delta \leftarrow \mathrm{DriftAttributor}(X_t, \mathcal{B}, A)$
  \STATE $\mathrm{attribution} \leftarrow \mathrm{TopK}(\Delta)$
  \STATE $\mathrm{UpdateBuffer}(\mathcal{B}, E_A)$
  \STATE $\mathrm{FineTune}(A, X_t)$
\ENDIF
\RETURN drift, novelty, attribution
\end{algorithmic}
\end{algorithm}

% ── 4. Experiments ────────────────────────────────────────────────────────────
\section{Experiments}
\label{sec:experiments}

\subsection{Dataset: NSL-KDD}

We evaluate on the NSL-KDD dataset~\cite{tavallaee2009nsl},
a widely used benchmark for network intrusion detection.
It contains 125,973 labelled connection records with 41 raw features
(numeric and categorical) encoded to 122 binary/normalised features
after one-hot encoding and min-max normalisation.
The dataset covers 22 distinct attack types across four categories:
DoS, Probe, R2L, and U2R.

\subsection{Experimental Protocol}

We simulate a realistic, non-stationary stream in three phases:
\begin{itemize}
  \item \textbf{Phase 1} (chunks 1--5): Normal traffic only.
        Vigil is trained offline on 1,000 normal samples.
  \item \textbf{Phase 2} (chunks 6--20): Known attack classes injected
        at 50\% concentration per chunk (100 normal + 100 attack).
  \item \textbf{Phase 3} (chunks 21--50): Novel attack class added at
        25\% concentration alongside Phase 2 attacks.
\end{itemize}

Each chunk contains 200 samples ($n=200$, $d=122$).
The stream totals 10,000 samples across 50 chunks.
Ground truth: chunks 1--5 are ``no drift''; chunks 6--50 are ``drift''.

\textbf{Vigil configuration}: buffer size = 600, drift threshold $\delta =
0.10$, $r = 15$ replications, $s = 30$ sample size, $k = 5$ top features,
400 training epochs.

\textbf{Baselines}: We compare against three classical detectors from the
\texttt{river} library~\cite{montiel2021river}:
ADWIN~\cite{bifet2007adwin},
KSWIN~\cite{raab2020kswin}, and
Page-Hinkley~\cite{page1954continuous}.
Classical detectors are fed a scalar error signal from a 1-nearest-neighbour
classifier trained on normal traffic, following the standard
univariate evaluation protocol.

\subsection{Drift Detection Results}

Table~\ref{tab:results} presents precision, recall, F1, and
detection delay across all methods.

\begin{table}[!t]
\caption{Drift Detection on NSL-KDD (50 chunks $\times$ 200 samples).
  ``Delay'' is chunks after drift injection until first detection.
  ``Attr.'' indicates feature-level attribution support.}
\label{tab:results}
\centering
\begin{tabular}{lcccccc}
\toprule
\textbf{Method} & \textbf{Prec.} & \textbf{Rec.} & \textbf{F1}
  & \textbf{Delay} & \textbf{Labels?} & \textbf{Attr.?} \\
\midrule
\textbf{Vigil (ours)} & \textbf{93.3\%} & \textbf{31.1\%}
  & \textbf{46.7\%} & \textbf{1 chunk} & No & \textbf{Yes} \\
ADWIN         &  0.0\% &  0.0\% &  0.0\% & N/A     & No & No \\
KSWIN         & 100.0\% & 2.2\% &  4.3\% & 8 chunks & No & No \\
Page-Hinkley  &  0.0\% &  0.0\% &  0.0\% & N/A     & No & No \\
\bottomrule
\end{tabular}
\end{table}

Vigil achieves \textbf{F1 = 46.7\%} with only 1 false positive out of 5
stable chunks (precision = 93.3\%). Classical univariate detectors
essentially fail: ADWIN and Page-Hinkley fire zero detections;
KSWIN fires once (chunk 13, 8 chunks after drift injection)
with recall of only 2.2\%.

The fundamental reason for this disparity is dimensionality:
the 1-NN error signal averaged across a mixed chunk (50\% normal
+ 50\% attack) produces a weak signal that falls within the normal
fluctuation range. Vigil's dual-autoencoder operates directly in the
full 122-dimensional space, detecting the increased reconstruction
difficulty caused by out-of-distribution samples.

Vigil's recall (31.1\%) reflects the inherent challenge of detecting
drift from mixed chunks where attack samples may share reconstruction-error
characteristics with normal samples (particularly subtle attacks such as
\texttt{ipsweep} and \texttt{loadmodule}).
This is a fundamental limitation of unsupervised, label-free detection
rather than a failure of the approach.

\subsection{Novel-Class Detection}

In Phase 3 (chunks 21--50), novel attack classes appear at 25\%
concentration per chunk. With the data-driven KDE threshold (5th percentile
of training density), Vigil achieves \textbf{100\% novelty recall}:
every novel-class chunk is flagged with a reported novelty proportion
averaging 30\%.

Critically, the KDE threshold is calibrated automatically from the
training data, requiring no manual tuning. The original fixed threshold
(0.02) from the base framework fails completely on this dataset due to
KDE density scale mismatch---a bug we identify and fix in our implementation.

\subsection{Feature Attribution Analysis}

Fig.~\ref{fig:top_features} shows the top-5 features ranked by
DriftAttributor across all 14 detected drift events.
\texttt{root\_shell} (14.7\%) and \texttt{service\_telnet} (9.3\%)
dominate, consistent with privilege escalation and remote service
exploitation patterns in DoS and R2L attacks.
The feature \texttt{service\_ecr\_i} (6.3\%) corresponds to ICMP echo
traffic characteristic of ping-flood DoS variants.

\begin{figure}[!t]
  \centering
  \includegraphics[width=0.95\columnwidth]{figures/fig1_top_features.pdf}
  \caption{Top-5 network features ranked by DriftAttributor contribution
           across all drift-detected chunks in the NSL-KDD experiment.}
  \label{fig:top_features}
\end{figure}

\begin{figure}[!t]
  \centering
  \includegraphics[width=0.95\columnwidth]{figures/fig2_comparison.pdf}
  \caption{Precision, Recall, and F1 comparison of Vigil vs.\ classical
           drift detectors on NSL-KDD (50-chunk stream).}
  \label{fig:comparison}
\end{figure}

\begin{figure}[!t]
  \centering
  \includegraphics[width=0.95\columnwidth]{figures/fig3_per_class.pdf}
  \caption{Per-attack-class detection rate. Subtle attacks
           (\texttt{ipsweep}, \texttt{loadmodule}) have near-zero
           reconstruction error delta, remaining below detection threshold.}
  \label{fig:perclass}
\end{figure}

Table~\ref{tab:perclass} provides per-attack-class detection rates.
Buffer-overflow (100\%) and back/DoS (60\%) are detected reliably;
port-scan variants (ipsweep, nmap) produce smaller distributional shifts
and are more often missed. These results directly inform which attack
families warrant separate drift-threshold tuning in deployment.

\begin{table}[!t]
\caption{Per-Attack-Class Detection Rate (Vigil, NSL-KDD)}
\label{tab:perclass}
\centering
\begin{tabular}{lcc}
\toprule
\textbf{Attack Class} & \textbf{Detection Rate} & \textbf{Avg. Severity} \\
\midrule
buffer\_overflow   & 100.0\% & 0.20 \\
back (DoS)         &  60.0\% & 0.47 \\
land               &  50.0\% & 0.37 \\
multihop           &  40.0\% & 0.16 \\
ftp\_write         &  25.0\% & 0.03 \\
nmap               &  25.0\% & 0.05 \\
ipsweep            &   0.0\% & 0.00 \\
loadmodule         &   0.0\% & 0.02 \\
\bottomrule
\end{tabular}
\end{table}

\subsection{Ablation: DriftAttributor Overhead}

DriftAttributor adds negligible computational overhead:
the per-feature delta computation is $\mathcal{O}(n \cdot d)$ per chunk,
completing in under 1\,ms on CPU for $n=200$, $d=122$.
Detection latency is dominated by the autoencoder forward pass and
replicated T-tests, which together complete in $<$50\,ms per chunk.

% ── 5. Production Deployment ──────────────────────────────────────────────────
\section{Production Deployment}

Vigil ships as a complete MLOps system beyond the core library:
\begin{itemize}
  \item \textbf{Kafka pipeline}: A consumer reads feature chunks from a
        Kafka topic, invokes Vigil, and publishes drift alerts to a
        downstream topic.
  \item \textbf{FastAPI service}: \texttt{/fit} and \texttt{/detect}
        endpoints enable deployment as a microservice.
  \item \textbf{Airflow DAG}: An auto-retrain workflow is triggered when
        drift accumulates beyond a configurable threshold, with a
        quality gate that prevents degraded models from being deployed.
  \item \textbf{MLflow tracking}: Every chunk's drift severity, novelty
        proportion, and attribution report are logged as run metrics.
  \item \textbf{Streamlit dashboard}: A SOC-facing live dashboard
        visualises drift severity, novelty rate, and top drifted
        features in real time.
\end{itemize}

% ── 6. Discussion ─────────────────────────────────────────────────────────────
\section{Discussion}

\textbf{Precision vs. Recall trade-off.}
Vigil's 93.3\% precision with 31.1\% recall reflects a deliberate
design for low false-alarm rate in SOC environments, where alert fatigue
is a primary operational concern. The drift threshold $\delta$ can be
lowered to increase recall at the cost of precision, providing a
tuneable operating point.

\textbf{Limitations.}
(1)~Recall is constrained by the mixed-chunk evaluation protocol;
production deployments with higher attack concentrations will see
higher recall. (2)~The current implementation processes chunks in a
single thread; distributed processing across Kafka partitions would
increase throughput. (3)~DriftAttributor ranks features by
reconstruction error delta, which measures correlation with the
autoencoder's latent representation rather than causal importance
in a strict sense.

\textbf{Future work.}
Causal attribution methods (e.g., SHAP over the autoencoder's
feature embeddings), per-class drift thresholds, and evaluation on
CICIDS2017~\cite{sharafaldin2018cicids} are natural extensions.

% ── 7. Conclusion ─────────────────────────────────────────────────────────────
\section{Conclusion}

We presented \textbf{Vigil}, an open-source, pip-installable drift
detection library for network traffic streams that closes the
explainability gap in unsupervised drift monitoring.
Our primary contribution, \textbf{DriftAttributor}, ranks input features
by their per-sample reconstruction error increase at the time of drift,
providing actionable intelligence to SOC analysts without any labels.

On NSL-KDD, Vigil achieves 93.3\% drift detection precision,
100\% novel-class recall, and 1-chunk detection delay---dramatically
outperforming classical univariate detectors that achieve near-zero recall
on the same benchmark. We also identify and fix a critical KDE
threshold bug in the base framework that previously prevented
novel-class detection entirely.

Vigil is released under the MIT licence and is available on PyPI
(\texttt{pip install vigil-drift}) with complete documentation,
Kafka/Airflow/MLflow integrations, and a Streamlit SOC dashboard.

% ── Acknowledgements ──────────────────────────────────────────────────────────
\section*{Acknowledgements}
The author thanks the open-source contributors of PyTorch, river,
scikit-learn, and Apache Kafka, whose tools made this work possible.

% ── References ────────────────────────────────────────────────────────────────
\bibliographystyle{IEEEtran}
\bibliography{vigil_paper}

\end{document}
